\subsection{Test functions}

In this issue, U denotes an open subset of \(\mathbf{R}^{n}\) . For \(p \in [1, +\infty]\) we will write \(\mathcal{L}^{p}(U)\) and \(L^{p}(U)\) rather than \(\mathcal{L}^{p}(U, \mu)\) and \(L^{p}(U, \mu)\) . We identify the continuous functions belonging to \(\mathcal{L}^{p}(U)\) with their images in \(L^{p}(U)\) .

Equip the space \(\mathcal{C}^{\infty}(\mathrm{U})\) of infinitely differentiable functions in U with complex values, with the topology defined by the family of seminorms \(p_{\alpha,K}\) defined for \(\alpha \in \mathbf{N}^{n}\) and \(K \subset U\) by

\[
p _ {\alpha , \mathrm{K}} (\varphi) = \sup _ {x \in \mathrm{K}} | \partial^ {\alpha} \varphi (x) |.
\]

This space is complete (cf. EVT, III, p. 9, example b)).

For every compact subset K of U, we denote by \(\mathcal{C}_{\mathrm{K}}^{\infty}(\mathrm{U})\) the subspace of \(\mathcal{C}^{\infty}(\mathrm{U})\) formed by functions with support in K. The space \(\mathcal{C}_{\mathrm{K}}^{\infty}(\mathrm{U})\) is equipped with the topology induced by that of \(\mathcal{C}^{\infty}(\mathrm{U})\) . It is a Fréchet space, a countable family of seminorms defining its topology being the family \((p_{\alpha,\mathrm{K}})_{\alpha\in\mathbf{N}^{n}}\) . In particular, it is a bornological space (EVT, III, p. 12, prop. 2).

We denote by \(\mathcal{D}(\mathrm{U})\) the vector space \(\mathcal{K}(\mathrm{U})\cap \mathcal{C}^{\infty}(\mathrm{U})\) of functions of class \(\mathrm{C}^{\infty}\) with compact support in U. We say that \(\mathcal{D}(\mathrm{U})\) is the space of test functions in U. The space \(\mathcal{D}(\mathrm{U})\) is the inductive limit of the spaces \(\mathcal{C}_{\mathrm{K}}^{\infty}(\mathrm{U})\) and it is equipped with the corresponding locally convex inductive limit topology (EVT, II, p. 31).

This space is denoted \(\mathcal{C}_{\circ}^{\infty}(U)\) in EVT, III, p. 9.

Let \((\mathrm{K}_{m})\) be an increasing sequence of compact subsets of U whose interiors form a covering of U. The space \(\mathcal{D}(\mathrm{U})\) is then the strict inductive limit of the spaces \(\mathcal{C}_{\mathrm{K}_{m}}^{\infty}(\mathrm{U})\) (EVT, III, p. 9). It is therefore a complete space (EVT, II, p. 35, prop. 9). Every bounded subset of \(\mathcal{D}(\mathrm{U})\) is contained in one of the subspaces \(\mathcal{C}_{\mathrm{K}_{m}}^{\infty}(\mathrm{U})\) (EVT, III, p. 5, prop. 6). This means that \(\mathrm{B}\subset\mathcal{D}(\mathrm{U})\) is bounded if and only if there exists a compact subset K of U and a family \((\mathrm{M}_{\alpha})_{\alpha\in\mathbf{N}^{n}}\) in \(R_{+}\) such that B is contained in the set of functions \(\varphi\in\mathcal{C}_{\mathrm{K}}^{\infty}(\mathrm{U})\) satisfying

\[
p _ {\alpha , \mathrm{K}} (\varphi) \leqslant \mathrm{M} _ {\alpha}
\]

for all \(\alpha\in\mathbf{N}^{n}\) .

The space \(\mathcal{D}(\mathrm{U})\) is a bornological space (EVT, III, p. 12, example 3) and a Montel space (EVT, IV, p. 18, example 4).

Let V be an open subset of \(\mathbf{R}^{n}\) contained in U. If \(K \subset V\) is compact, then the restriction of functions to V defines a continuous and surjective linear map from \(\mathcal{C}_{\mathrm{K}}^{\infty}(\mathrm{U})\) onto \(\mathcal{C}_{\mathrm{K}}^{\infty}(\mathrm{V})\) . The extension by zero of a function defined on V induces an injective continuous linear map, called canonical, from \(\mathcal{D}(\mathrm{V})\) into \(\mathcal{D}(\mathrm{U})\) .

Remark. --- One defines similarly the space \(\mathcal{D}_{\mathbf{R}}(\mathrm{U})\) of test functions in U with real values. The linear map such that \(z \otimes \varphi \mapsto z\varphi\) for all \(z \in \mathbf{C}\) and every \(\varphi \in \mathcal{D}_{\mathbf{R}}(\mathrm{U})\) is an isomorphism of \(\mathbf{C} \otimes \mathcal{D}_{\mathbf{R}}(\mathrm{U})\) into \(\mathcal{D}(\mathrm{U})\) .

Lemma 2. --- Let \(\mathcal{U}\) be an open covering of U. There exists a locally finite open covering of U finer than \(\mathcal{U}\) and consisting of relatively compact sets.

Since U is locally compact, there exists a covering \(\mathcal{V}\) of U finer than \(\mathcal{U}\) consisting of open sets relatively compact in U. Since U is paracompact (TG, IX, p. 51, th. 4), there exists a locally finite open covering \(\mathcal{W}\) finer than \(\mathcal{V}\). Then \(\mathcal{W}\) is finer than \(\mathcal{U}\) and consists of relatively compact open sets.

Lemma 3. --- Let \(f \in \mathcal{K}(\mathbf{R}^n)\) and let V be an open neighborhood of the support K of f. There exists an integer \(m_0 \geqslant 1\) such that, for every \(x \in K\) , the set V contains the ball of radius \(m_0^{-1}\) centered at x. For every integer \(m > m_0\) , let \(V_m\) be the closed ball of radius \(m^{-1}\) centered at 0 in \(\mathbf{R}^{n}\) and let \(\varphi_{m}\) be a positive test function with support in \(V_{m}\) and integral 1. Set \(f_{m} = \varphi_{m} * f\) .

\begin{enumerate}
\def\labelenumi{\alph{enumi})}
\item
  One has \(f_{m}\in \mathcal{D}(\mathbf{R}^{n})\) and the support of \(f_{m}\) is contained in V;
\item
  Let \(p \in [1, +\infty]\) . The sequence \((f_m)_{m > m_0}\) converges to \(f\) in \(L^p(\mathbf{R}^n)\) .
\end{enumerate}

According to TG, IX, p. 14, remark, we have \(d(\mathrm{K}, \mathbf{R}^{n}-\mathrm{V}) > 0\) . Let \(m_{0} \geqslant 1\) be an integer such that \(m_{0}^{-1} < d(\mathrm{K}, \mathbf{R}^{n}-\mathrm{V})\) ; it satisfies the required condition. For all \(m > m_{0}\) , the function \(f_{m}\) is continuous (INT, VIII, p. 166, § 4, n° 5, prop. 11) and has support contained in \(\mathrm{K} + \mathrm{V}_{m}\) (INT, VIII, p. 126, § 1, n° 4, prop. 5), hence in V. By Corollary 2 of IV, p. 198, we have \(f_{m} \in \mathcal{D}(\mathbf{R}^{n})\) . If \(p \neq +\infty\) , then the sequence ( \(f_{m}\) ) converges to \(f\) in \(\mathrm{L}^{p}(\mathbf{R}^{n})\) (INT, VIII, p. 172, § 4, no. 7, prop. 20). \(^{(1)}\)

Suppose that \(p = +\infty\) . Let \(\varepsilon > 0\) . The function \(f\) is uniformly continuous on \(\mathbf{R}^n\) ; therefore there exists an integer \(m_1 > m_0\) such that, for every \(m \geqslant m_1\) and for all \(x \in \mathbf{R}^n\) and \(y \in \mathrm{V}_m\) , one has \(|f(x) - f(x - y)| < \varepsilon\) .

For every \(m \geqslant m_{1}\) , the function \(\varphi_{m}\) vanishes outside \(\mathrm{V}_{m}\) , is positive, and has integral 1. We therefore deduce the inequality

\[
| f (x) - f _ {m} (x) | = \left| \int_ {\mathbf {R} ^ {n}} (f (x) - f (x - y)) \varphi_ {m} (y) d \mu (y) \right| \leqslant \varepsilon
\]

for all \(x \in \mathbf{R}^{n}\) , whence the result.

Proposition 4. --- a) The canonical injection of \(\mathcal{D}(\mathrm{U})\) into \(\mathcal{K}(\mathrm{U})\) is continuous and \(\mathcal{D}(\mathrm{U})\) is dense in \(\mathcal{K}(\mathrm{U})\) ;

\begin{enumerate}
\def\labelenumi{\alph{enumi})}
\setcounter{enumi}{1}
\tightlist
\item
  For every \(p \in [1, +\infty[\) , the canonical injection of \(\mathcal{D}(\mathrm{U})\) into \(\mathcal{L}^p(\mathrm{U})\) is continuous and \(\mathcal{D}(\mathrm{U})\) is dense in \(\mathrm{L}^p(\mathrm{U})\) .
\end{enumerate}

Every continuous seminorm on \(\mathcal{K}(\mathrm{U})\) is continuous on \(\mathcal{D}(\mathrm{U})\) , hence the canonical injection of \(\mathcal{D}(\mathrm{U})\) into \(\mathcal{K}(\mathrm{U})\) is continuous.

Let \(f \in \mathcal{K}(\mathrm{U})\) . There exists a sequence \((f_{m})_{m \in \mathbf{N}}\) in \(\mathcal{D}(\mathrm{U})\) which converges uniformly to f on U and such that the supports of the \(f_{m}\) are contained in a fixed relatively compact neighborhood of the support of f (Lemma 3). The sequence \((f_{m})\) therefore converges to f in \(\mathcal{K}(\mathrm{U})\) . This concludes the proof of a).

Let \(p \in [1, +\infty[\) . Let \(f \in \mathcal{D}(U)\) and let K be the support of f. We then have the inequality \(\mathrm{N}_{p}(f) \leqslant \mu(\mathrm{K})^{1/p} \sup_{x \in \mathrm{K}} |f(x)|\) , hence the canonical injection of \(\mathcal{D}(U)\) into \(\mathcal{L}^{p}(U)\) is continuous. The last part of assertion b) then follows from a).

The product of two test functions being again a test function, the Leibniz formula (FVR, I, p. 28, prop. 2) shows that the space \(\mathcal{D}(\mathrm{U})\) is a topological algebra.

More generally, let \(f \in \mathcal{C}^{\infty}(\mathrm{U})\) . The linear map \(\varphi \mapsto f\varphi\) from \(\mathcal{D}(\mathrm{U})\) into \(\mathcal{D}(\mathrm{U})\) is then continuous. Indeed, for every compact subset \(\mathbf{K}\) of \(\mathrm{U}\) and every \(\alpha \in \mathbf{N}^{n}\) , we have

\[
p _ {\alpha , \mathrm{K}} (f \varphi) \leqslant \sum_ {0 \leqslant \beta \leqslant \alpha} \binom {\alpha} {\beta} p _ {\beta , \mathrm{K}} (f) p _ {\alpha - \beta , \mathrm{K}} (\varphi)
\]

(cf. loc. cit.).

